% ===========================================================================
% Conclusion générale.
% Aucune preuve nouvelle ici : uniquement le bilan de ce qui a été démontré.
% ===========================================================================
\chapter*{Conclusion générale}
\addcontentsline{toc}{chapter}{Conclusion générale}

Ce projet de fin d'études a porté sur la conception et la réalisation de
PIQBIT, une plateforme de recrutement augmentée par l'apprentissage
automatique. Le problème traité était celui d'une double contrainte : traiter
automatiquement un volume important de candidatures, tout en produisant des
décisions explicables et contestables par un recruteur humain.

\section*{Bilan des objectifs}

Le tableau~\ref{tab:objectives_review} reprend chaque objectif spécifique du
chapitre~\ref{chap:context} et son état au terme du projet.

\begin{table}[htbp]
  \centering
  \caption{Bilan des objectifs spécifiques}
  \label{tab:objectives_review}
  \begin{tabularx}{\textwidth}{L{0.8cm} L{4.4cm} L{2.6cm} X}
    \toprule
    \textbf{N°} & \textbf{Objectif} & \textbf{État} & \textbf{Réserve} \\
    \midrule
    1 & Socle sécurisé            & Atteint & Écart E-03 sur un endpoint à corriger \\
    2 & Offres et candidatures    & Atteint & — \\
    3 & Analyse et notation explicable & Atteint & — \\
    4 & Rapprochement sémantique  & Atteint & Seul le repli est validé \\
    5 & Étapes aval du recrutement& Atteint & — \\
    6 & Modèle de prédiction salariale & Atteint & Métriques établies sur le jeu disponible \\
    7 & Validation et portes de qualité & \textbf{Partiellement atteint} & Porte fermée au 22 septembre 2026 \\
    \bottomrule
  \end{tabularx}
\end{table}

Six objectifs sur sept sont atteints. Le septième ne l'est que partiellement :
la chaîne de validation existe et s'exécute, mais sa porte de qualité était
fermée à la date de la campagne, trois tests étant en échec. Cette réserve est
énoncée plutôt qu'arrondie : une chaîne qui échoue est une chaîne qui
fonctionne — elle signale un défaut — mais elle ne permet pas de déclarer
l'objectif tenu.

\section*{Contributions d'ingénierie}

Quatre contributions se démarquent, en ce qu'elles ont exigé un arbitrage
plutôt qu'une mise en œuvre.

\paragraph{Une notation explicable par construction.} Le score n'est pas un
nombre opaque : il est décomposé en sous-scores persistés, rattachés aux
exigences de l'offre. Le choix de \emph{stocker} la décomposition plutôt que de
la recalculer garantit que ce qui a été montré à un recruteur reste consultable
ensuite, même si la logique de notation évolue. C'est cette propriété qui rend
une décision contestable — donc défendable.

\paragraph{La correction d'une surnotation silencieuse.} Le repli de
normalisation des compétences, appliqué globalement, faisait disparaître les
compétences non reconnues dès qu'une seule l'était, et surnotait les candidats
sans produire la moindre erreur. Identifier un défaut qui ne se manifeste par
aucun symptôme, puis construire le cas de test qui le révèle, a constitué
l'exercice de diagnostic le plus formateur du projet.

\paragraph{Une méthodologie de modélisation défendable.} Le modèle salarial
n'a pas été retenu parce qu'il donnait de bons chiffres, mais parce que ses
chiffres ont été mesurés correctement : conversion net-brut par recherche
dichotomique sur le barème réel plutôt que par coefficient forfaitaire,
sélection sur l'erreur de test plutôt que d'entraînement, retrait d'une
variable colinéaire identifiée comme mode de défaillance, et exclusion unique
conservée avec son motif écrit pour rester auditable.

\paragraph{Des dépendances externes toutes facultatives.} Messagerie, agenda,
visioconférence, modèles de langage : aucune n'est sur le chemin critique, et
chacune possède un repli implémenté. La plateforme reste utilisable, en service
dégradé et annoncé, lorsque l'ensemble est indisponible.

\section*{Limites}

Les douze limites recensées au chapitre~\ref{chap:validation} ne se répètent
pas ici, mais trois méritent d'être rappelées parce qu'elles conditionnent tout
usage réel de la plateforme.

L'écart d'autorisation E-03 doit être corrigé avant toute mise en service :
un rôle non habilité atteint une action qui engage financièrement
l'entreprise. L'absence d'environnement cible et de chiffrement du transport
interdit un déploiement en l'état. Enfin, l'absence totale de télémétrie
signifie qu'une dégradation du service ne serait détectée que par ses
utilisateurs.

Ces limites ne remettent pas en cause la validité de la démarche : elles
délimitent ce qui a été démontré de ce qui reste à établir.

\section*{Perspectives}

Les extensions les plus utiles se situent à trois niveaux.

\emph{À court terme}, la fermeture des écarts : correction de E-03, réécriture
des deux tests périmés, nouvelle campagne datée, réduction du bundle sous son
budget, et capture d'une exécution du chemin vectoriel natif.

\emph{À moyen terme}, l'industrialisation : terminaison TLS, environnement
cible avec procédure de retour arrière, télémétrie minimale sur les parcours
critiques, analyse de dépendances et recherche de secrets dans les portes de
qualité.

\emph{À plus long terme}, l'enrichissement fonctionnel : extension de l'analyse
à d'autres langues, mesure et correction des biais de la notation — un système
qui classe des personnes doit pouvoir démontrer qu'il ne reproduit pas des
écarts systématiques — et boucle de retour permettant d'ajuster la pondération
des sous-scores à partir des recrutements effectivement aboutis.

Ce dernier point est peut-être le plus important. PIQBIT établit aujourd'hui
une correspondance entre un CV et une offre ; il ne sait pas encore si cette
correspondance prédit une embauche réussie. Relier la notation aux résultats
observés transformerait un outil de tri en outil d'apprentissage.
